\section{Level C: reputation}\label{sec:levelc}

Reputation enters in two separate scores on unlinkable pseudonyms (\cref{req:separate}). The
author score limits how much an author may propose. The evaluator score weights reviews.

\subsection{Author score}

Let $q_j\in[0,1]$ be the final quality of author $a$'s item $j$, a function of its Level~B
statistics whose exact definition is still open. The author score is the shrunken, time-discounted
average
\begin{equation}\label{eq:author}
  C_a=\frac{\alpha_0+\sum_j w_jq_j}{\alpha_0+\beta_0+\sum_j w_j},\qquad w_j=e^{-\Delta t_j/T},
\end{equation}
with prior $(\alpha_0,\beta_0)=(2,3)$ and $T=18$ months. Two accepted items out of two give
$4/7\approx0.57$; 180 out of 200 give $182/205\approx0.89$. $C_a$ sets the proposal rate
$q_{\min}+(q_{\max}-q_{\min})C_a$. It never weights a vote \designchoice. A failed appeal to
evidence is to be recorded as a negative pseudo-observation inside $C_a$, escrowed when the appeal
is filed (design decision D27). The reference implementation still applies a simpler
stake-and-refund rule.

\subsection{Evaluator score: forecasts scored on golden items}

Every reviewer declares a probability $p_{uj}$ that item $j$ will pass empirical validation. The
same number is the rating $r_{uj}$ of Level~A. About 5\% of every review queue consists of
\emph{golden} items. These are indistinguishable from the rest, and their outcome
$o_j\in\{0,1\}$ is known from validated history (design decision D21). The forecasts on golden
items are scored. The design goal (design decision D6) is that honest forecasting is optimal, that
copying the crowd earns nothing, and that being right where the crowd is wrong is what pays. Its
opposite would be a Schelling-point mechanism that rewards agreement with the majority. The
current rule is the Brier skill score against the crowd~\cite{brier1950,murphy1988skill}:
\begin{equation}\label{eq:bss}
  \mathrm{BSS}_u=1-\frac{\sum_j(p_{uj}-o_j)^2}{\sum_j(\bar p_j-o_j)^2},\qquad
  \bar p_j=\frac{\sum_v w_vp_{vj}}{\sum_v w_v},
\end{equation}
where the crowd forecast $\bar p_j$ averages the whole panel, the reviewer included (design
decision D23). The evaluator score is $E_u=\sig(\gamma\,\mathrm{BSS}_u)$. It is updated by an
asymmetric moving average that rises slowly and falls fast, and it gives the weight
$w_u=\min(w_{\max},E_u)$ with $w_{\max}=3\operatorname{median}(w)$.

Scoring only golden items is more than a convenience. On a live item the outcome is observed only
if the item reaches Level~B, and whether it does depends on the forecasts themselves, because they
are also the bridging ratings. When a forecast influences whether its outcome is observed, known
results on decision scoring show that strict properness requires randomizing the
decision~\cite{othman2010decision,chen2014eliciting}.
Golden items avoid the problem because their outcome is fixed in advance.

\subsection{The ratio-form skill score is not proper}

\begin{proposition}[Non-properness of the skill score]\label{prop:bss}\proved{}
Score a single item with baseline $b\in(0,1)$ that does not depend on the report, and let the
reviewer believe $P(o=1)=q\in(0,1)$. The expected skill score
\[
  \E[\mathrm{BSS}]=1-\frac{q(1-p)^2}{(1-b)^2}-\frac{(1-q)\,p^2}{b^2}
\]
is strictly concave in $p$, and its maximizer satisfies
\begin{equation}\label{eq:bss-opt}
  \logit p^{*}=\logit q+2\logit b .
\end{equation}
Truthful reporting is optimal if and only if $b=\tfrac12$.
\end{proposition}
\begin{proof}
The expectation over $o$ is immediate. The first-order condition
$q(1-p)/(1-b)^2=(1-q)p/b^2$ gives $p/(1-p)=\frac{q}{1-q}\cdot\frac{b^2}{(1-b)^2}$, which is
\eqref{eq:bss-opt}. The second derivative is $-2q/(1-b)^2-2(1-q)/b^2<0$.
\end{proof}

That skill scores built on a proper score are themselves improper, and that the resulting hedging
shrinks as the number of scored forecasts grows, is known in forecast
verification~\cite{murphy1973hedging}. What matters here is the direction of the hedge. A ratio of
two random sums is not an expectation of a score. The denominator gives more weight to outcomes
where the crowd was right, and \eqref{eq:bss-opt} shows the result: the report is pushed
toward the outcome the crowd favours, by twice the crowd's log-odds. With the crowd at $0.65$,
a reviewer who believes $0.78$ should report $0.92$, and a reviewer who believes $0.30$ should
report $0.60$, crossing to the crowd's side (\cref{fig:bss}). The distortion is largest for the
dissenter, who is exactly the reviewer the design wants to reward.

\begin{figure}[t]
\centering
\begin{tikzpicture}
\begin{axis}[
  width=0.62\linewidth, height=5.6cm,
  xlabel={belief $q$}, ylabel={optimal report $p^*$},
  xmin=0, xmax=1, ymin=0, ymax=1,
  legend style={font=\scriptsize, at={(0.98,0.03)}, anchor=south east, draw=none, fill=white, fill opacity=0.8, text opacity=1},
  grid=major, grid style={gray!20}, tick label style={font=\scriptsize}, label style={font=\small}]
\addplot[gray, dashed, domain=0:1] {x};
\addlegendentry{truthful}
\addplot[thick, blue!70!black] table[col sep=comma, x=q, y=pstar_fixed] {data/levelC_bss_curve.csv};
\addlegendentry{BSS, baseline $0.65$ (closed form)}
\addplot[thick, orange!85!black, densely dotted] table[col sep=comma, x=q, y=pstar_selfincl] {data/levelC_bss_curve.csv};
\addlegendentry{BSS, baseline includes the reviewer}
\addplot[only marks, mark=*, mark size=1.8pt, red!70!black] coordinates {(0.30,0.5965)};
\end{axis}
\end{tikzpicture}
\caption{Report that maximizes the expected Brier skill score on one golden item, against the
reviewer's belief, when the other panelists' mean forecast is $0.65$. Solid line: baseline
excluding the reviewer (\cref{prop:bss}). Dotted line: baseline including the reviewer with
weight $1/9$, as in the implementation (numerical). The marked point is a dissenter who
believes $0.30$ and is best off reporting $0.60$. Source: \texttt{scripts/levelC\_bss.py}.}
\label{fig:bss}
\end{figure}

With several items the ratio couples them, and there is no closed form. \Cref{tab:bss} gives the
exact optimum, enumerating all $2^m$ outcomes, for random beliefs and baselines.

\begin{table}[t]
\centering\small
\caption{Distortion of the optimal report under the ratio-form skill score, $|p^*_j-q_j|$,
for $m$ jointly scored items with beliefs uniform on $[0.1,0.9]$ and baselines uniform on
$[0.2,0.8]$ (200 random instances; 60 for $m=8$). Source: \texttt{scripts/levelC\_bss.py}.}
\label{tab:bss}
\begin{tabular}{@{}rcccc@{}}
\toprule
& \multicolumn{2}{c}{baseline: others only} & \multicolumn{2}{c}{baseline includes reviewer}\\
\cmidrule(lr){2-3}\cmidrule(l){4-5}
$m$ & mean & max & mean & max\\
\midrule
1 & 0.229 & 0.577 & 0.227 & 0.569\\
2 & 0.127 & 0.434 & 0.118 & 0.416\\
4 & 0.061 & 0.173 & 0.061 & 0.227\\
8 & 0.030 & 0.081 & 0.035 & 0.095\\
\bottomrule
\end{tabular}
\end{table}

\begin{finding}\label{find:bss}\empirical{}
The mean distortion falls roughly as $0.23/m$ in the number of jointly scored golden items,
consistent with Murphy's observation. Including the reviewer in the baseline, as the
implementation does, changes it little. At a golden-item rate of 5\%, a reviewer who judges about
ten items per epoch meets a golden item about every other epoch. A score computed over a few
golden items is therefore in the regime of the largest distortion.
\end{finding}

\subsection{A proper alternative}

\begin{proposition}[Leave-one-out difference score]\label{prop:loo}\proved{}
Let $\bar p_{-u,j}$ be the weighted mean forecast of the \emph{other} panelists on item $j$, and
\begin{equation}\label{eq:loo}
  S_u=\frac1m\sum_{j}\Bigl[(\bar p_{-u,j}-o_j)^2-(p_{uj}-o_j)^2\Bigr].
\end{equation}
Under \cref{ass:blind}, so that $\bar p_{-u,j}$ does not depend on $p_{uj}$:
\begin{enumerate}[label=(\roman*)]
  \item $S_u$ is strictly proper: $\E[S_u]=c_u-\frac1m\sum_j(p_{uj}-q_j)^2$ with $c_u$ not
    depending on the reports, so truthful reporting is the unique optimum;
  \item a reviewer who copies the crowd ($p_{uj}=\bar p_{-u,j}$) scores exactly $0$ for every
    outcome, not only in expectation;
  \item $S_u>0$ if and only if the reviewer's Brier score on the scored items beats the crowd's;
  \item $S_u$ is additive over items, so it aggregates across epochs without reweighting.
\end{enumerate}
\end{proposition}
\begin{proof}
For (i), use $\E[(p-o)^2]=(p-q)^2+q(1-q)$ for $o\sim\mathrm{Bernoulli}(q)$~\cite{gneiting2007scoring}.
Parts (ii)--(iv) are immediate from \eqref{eq:loo}.
\end{proof}

The difference score keeps what D6 wants from the skill score, namely zero for copying the crowd
and a positive score only for beating it, and drops the ratio that breaks properness. One new
exposure deserves mention. Because the baseline is formed by the other panelists, a coalition can
raise a member's score by degrading its own forecasts. This costs the saboteurs their own scores,
is bounded by their share of the panel weight, and is blunted by the fact that golden items are
indistinguishable from live ones, whose ratings also enter bridging.

\subsection{From scores to weights}

\begin{proposition}[A cap that never binds]\label{prop:cap}\proved{}
If all weights lie in $(0,1)$ and $\operatorname{median}(w)\ge\tfrac13$, then
$w_{\max}=3\operatorname{median}(w)\ge1>w_u$ for every $u$, so $\min(w_{\max},E_u)=E_u$ always.
\end{proposition}

With $E_u=\sig(\gamma\,\mathrm{BSS}_u)$ and a crowd baseline that puts the typical reviewer's
score near zero, the median of $E$ is close to $\tfrac12$, so the cap is vacuous in the regime the
design expects. The repository audit notes the same point. Expressing weights relative to the
median, $w_u=E_u/\operatorname{median}(E)$, or on the odds scale, $w_u=e^{\gamma S_u}$, would give
the cap force \openq.

A subtler issue concerns incentives. The weight is a non-linear function of the realized score,
first through $\sig$ and then through a moving average that falls faster than it rises (a guard
against the long con). A reviewer who maximizes expected future weight therefore faces a
loss-averse objective, and reducing variance pays. Under \cref{prop:loo} the zero-variance
strategy is to copy the crowd. The asymmetry that deters the long con thus pulls in the same
direction as herding, against D6. How strongly depends on $\gamma$, on the two rates and on the
number of scored items, none of which the design specifies \openq. A natural mitigation is to
accumulate the linear score \eqref{eq:loo} over a long window and map it to a weight only
afterwards, so that each item's contribution stays small and truthful reporting remains
approximately optimal.

\subsection{Judgments without ground truth}

Some dimensions never receive an empirical verdict, such as clarity or the quality of the cited
source. For these the design proposes peer prediction. Candidates include the output-agreement
score of Dasgupta and Ghosh~\cite{dasgupta2013crowdsourced}, which rewards agreement with a peer
on a shared item minus the agreement on unrelated items, multi-task correlated
agreement~\cite{shnayder2016informed}, and the Bayesian truth serum~\cite{prelec2004bts}. Only the
pairwise Dasgupta--Ghosh score is implemented. How it is aggregated over pairs and items, and the
incentive analysis in the presence of reputation weights, remain open \openq.
